% Generated draft: requires mathematical and source review.
\documentclass[11pt]{article}
\usepackage[margin=25mm]{geometry}
\usepackage{fontspec}
\setmainfont{Latin Modern Roman}
\usepackage{amsmath,amssymb,mathtools,unicode-math}
\setmathfont{Latin Modern Math}
\usepackage{enumitem,needspace,xurl,hyperref}
\setlength{\parindent}{0pt}
\setlength{\parskip}{5pt}
\setlength{\emergencystretch}{3em}
\newcommand{\N}{\mathbb{N}}
\newcommand{\Z}{\mathbb{Z}}
\newcommand{\Q}{\mathbb{Q}}
\newcommand{\R}{\mathbb{R}}
\newcommand{\C}{\mathbb{C}}
\newcommand{\F}{\mathbb{F}}
\newcommand{\A}{\mathbb{A}}
\newcommand{\PP}{\mathbb{P}}
\newcommand{\E}{\mathbb{E}}
\newcommand{\eps}{\varepsilon}
\newcommand{\dd}{\,\mathrm{d}}
\begin{document}
\section*{UnsolvedMath ID: 30005520}
\textbf{Source problem code:} catalogue entry
\subsection{Definitions and mathematical statement}
Phase retrieval concerns reconstructing an unknown object from measurements that do not directly provide its phase; tomographic reconstruction concerns reconstructing an object from indirect measurements acquired through an imaging geometry. These descriptions identify the subject only: the supplied record specifies neither an object space nor a measurement operator.

The source context is phase retrieval and tomographic reconstruction in X-ray near-field diffractive imaging. The supplied research question is: how much further improvement can be obtained by fully exploiting indirect measurements, sophisticated mathematical analysis, and high-performance algorithms, and what are the ultimate limits of such improvement?

This is a qualitative research question rather than a specified mathematical conjecture. The record does not identify the admissible objects, acquisition geometry, forward model, measurement budget, noise model, reconstruction criterion, computational resources, or baseline against which improvement is measured. It also does not specify whether the intended limits concern uniqueness, stability, resolution, statistical accuracy, computational efficiency, or a combination of these. Consequently, no definite quantified assertion or optimization problem can be extracted without additional hypotheses. These qualifications remain unresolved.
\subsection{Short English statement}
The question asks how far phase retrieval and X-ray near-field tomographic imaging can improve through better use of indirect data, mathematical analysis, and computation. A precise mathematical target requires a specified imaging model, constraints, and measure of improvement; these are absent from the supplied record.
\subsection{Sources}
\begin{itemize}
\item Supplied local UnsolvedMath registry record, permanent ID \(30005520\), source problem code \texttt{catalogue entry}, titled \emph{Limits of Phase Retrieval and Tomographic Imaging}, in the Oberwolfach Reports Open Problems collection.
\item Supplied report citation: \emph{Tomographic Inverse Problems: Mathematical Challenges and Novel Applications} (2023), Oberwolfach Reports, pp.\ 1105--1194, DOI \texttt{10.4171/owr/2023/21}. Source URL and supplied literature-source link, labeled ``Oberwolfach problem report'': \url{https://doi.org/10.4171/owr/2023/21}. No question-specific page is supplied.
\item The original extraction associates the question with \emph{Phase Retrieval and Tomographic Reconstruction in X-ray Near-field Diffractive Imaging: Inverse Problems at Work} and Tim Salditt. The registry gives proposed year 2023 but leaves the proposer field unspecified; the association does not establish explicit proposer attribution.
\item Literature assessment dated exactly \texttt{2026-08-22}: the supplied curation labels the question open and states that phase retrieval and tomographic algorithms already exploit indirect measurements substantially, but that no general theorem answers the question about ultimate improvement across mechanisms and noise regimes. This is supplied dated curation, not independently verified literature assessment; its status and cited claims require independent verification.
\item The registry labels the statement \texttt{corrected\_verified} and reports direct checking of the official report, claiming that the cleaned statement preserves the complete mathematical target while restoring definitions and removing extraction artifacts or source-dependent wording. This is a supplied verification claim, not verification performed for this draft.
\end{itemize}
\end{document}
